% Generated by code/08_generate_results_tables.R
\begin{table}[H]
\centering
\scriptsize
\caption[Reliability scorecard for the behavioural module]{Reliability scorecard for the behavioural module}
\label{tab:stage3_reliability_scorecard}
\resizebox{\textwidth}{!}{%
\begin{tabular}{@{}L{3.0cm}L{6.3cm}L{5.4cm}@{}}
\toprule
Diagnostic & Result & Weight for interpretation \\
\midrule
Preferred FE estimate & $\hat{\beta}_{\tau}=$ -1.500; semi-elasticity $=$ 1.50; $N=$ 8,971; pairs $=$ 1,936; within $R^2=$ 0.0339. & Moderate: the sign and magnitude are economically interpretable. \\
Model selection & BIC selects linear; $\Delta$BIC log $=$ 5.40 and quadratic $=$ 7.54; tax-term Wald $p=$ $1.5 \times 10^{-6}$. & Moderate: the public panel prefers a parsimonious linear response. \\
Clustered inference & Two-way clustered SE for ETR $=$ 0.272 ($p=$ $1.5 \times 10^{-6}$); affiliate-clustered SE $=$ 0.200 ($p=$ $5.6 \times 10^{-12}$). & Useful: two-way clustering is the main inference; affiliate clustering is a less conservative comparison. \\
FE-2SLS first stage & Simulated EATR: coefficient 0.247, $F=$ 12.05; lagged statutory CIT: coefficient 0.253, $F=$ 17.49. & Limited: the instruments are relevant enough for diagnostics. \\
Reduced forms and 2SLS & Reduced-form signs are negative but imprecise: simulated EATR -0.099 ($p=$ 0.829), lagged CIT -0.270 ($p=$ 0.515); FE-2SLS $p=$ 0.824 and 0.507. & Weak for causal replacement: reduced forms and 2SLS are too imprecise. \\
Low-ETR tail support & Below 5\% ETR: 17.4\% of rows and 10.1\% of profit; 5--15\% ETR: 27.3\% of rows. & Limited for curvature: the low-rate tail exists but does not strongly identify nonlinear extrapolation. \\
Leave-one affiliate & Affiliate omissions: semi-elasticity 1.45--1.51; model stable in 13 of 13; preferred revenue change -27.7\% to 2.7\%. & Moderate for the coefficient; weaker for final revenue because conduit jurisdictions are influential. \\
Leave-one parent & Parent omissions: semi-elasticity 1.40--1.60; model stable in 12 of 13; preferred revenue change -27.9\% to -0.1\%. & Moderate for the coefficient; model selection changes in one parent omission. \\
Final decision & Use the behavioural input as a scenario coefficient; do not make it the causal centrepiece. & The revenue decomposition should carry more interpretive weight than the behavioural module. \\
\bottomrule
\end{tabular}
}

\vspace{0.5em}
\begin{minipage}{0.96\textwidth}
\footnotesize\raggedright \textit{Notes:} The scorecard reports whether the behavioural module is usable for the revenue counterfactual. The preferred estimate is the linear two-way fixed-effect Hines-Rice coefficient with parent-and-affiliate clustered inference. The final row states the evidentiary decision used in the thesis. Source files: data/processed/diagnostics/step5\_stage3\_coefficients.csv; data/processed/diagnostics/step5\_stage3\_model\_comparison.csv; data/processed/diagnostics/step5\_stage3\_tax\_endogeneity\_diagnostics.csv; data/processed/diagnostics/step5\_stage3\_etr\_tail\_support.csv; data/processed/diagnostics/step10\_jurisdiction\_influence\_case\_summary.csv.
\end{minipage}
\end{table}
